\documentclass[12pt]{article}
\usepackage[utf8]{inputenc}
\usepackage[T1]{fontenc}
\usepackage{hyperref}
\title{Tboli Language Resources for K-12 Curriculum Development}
\author{Compiled by Claude Sonnet 5 (claude-sonnet-5) for Andr\'as Kornai}
\date{}
\begin{document}
\maketitle

\section{Introduction}
Tboli (ISO 639-3 tbl; Glottocode tbol1240) is an Austronesian language of the Philippines with roughly 95,000 speakers. It is written in the Latin script. Search turned up a dictionary and a basic word list. It turned up no primers, textbooks or MTB-MLE (mother tongue-based multilingual education) classroom materials. Those may exist but were not found in search results.

\section{Primary Educational Resources}
No Tboli primers, textbooks or teacher guides were found in the search results. The Tboli--English dictionary below is the closest thing to a school-usable resource. A follow-up search of Department of Education MTB-MLE materials and SIL Philippines literacy projects is recommended.

\section{Linguistic and Reference Materials}
\begin{itemize}
\item Tboli--English Dictionary (Awed, Silin, Moran and Underwood; SIL International, web edition 2018). The search snippet reports about 6,882 Tboli entries and 6,096 English entries, with semantic domains, grammar and phonology sections. \url{https://www.webonary.work/tboli/?p=146817}
\item Webonary Tboli site, further pages of the same dictionary. \url{https://www.webonary.work/tboli/?p=9444}
\item Webonary copyright and overview page. \url{https://www.webonary.org/tboli/overview/copyright/}
\item Tboli-English dictionary (SIL, Manila, 2004), library record at the University of the Philippines. \url{https://ds.mainlib.upd.edu.ph/Record/UP-99796217604438412}
\item Austronesian Basic Vocabulary Database, Tboli (Tagabili) word list. \url{https://abvd.eva.mpg.de/austronesian/language.php?id=133}
\end{itemize}

\section{Religious and Community Resources}
The Wikipedia entry reports that literature, a grammar, texts and the Bible exist in Tboli. No specific Bible title or URL for Tboli was found in the search results.
\begin{itemize}
\item Wikipedia, ISO 639 entry for Tboli (general background only). \url{https://wikipedia.com/wiki/ISO_639:tbl}
\end{itemize}

\section{Digital and Audio Resources}
The only digital resources found are the online Webonary dictionary and the ABVD word list listed above. No audio or video was found.

\section{Curriculum Development Context}
The reference description says Tboli is used as a language of instruction. The search found no curriculum documents to confirm this. The dictionary, with its semantic-domain organisation, could support vocabulary and literacy materials. Primers, readers and teacher materials would probably have to be developed or located through DepEd and SIL Philippines.
\end{document}